\documentclass[openany]{memoir}
\usepackage[utf8]{inputenc}
\usepackage[T1]{fontenc}

\usepackage{url,listings,courier,graphicx,bbding,xspace}

\newcommand{\gbeta}{\textsf{gbeta}\xspace}
\renewcommand{\beta}{\textsc{Beta}\xspace}
 \lstset{
         basicstyle=\normalsize\ttfamily, % Standardschrift
         numbers=left,               % Ort der Zeilennummern
         numberstyle=\tiny,          % Stil der Zeilennummern
         %stepnumber=2,               % Abstand zwischen den
         %Zeilennummern
         numbersep=5pt,              % Abstand der Nummern zum Text
         tabsize=2,                  % Groesse von Tabs
         extendedchars=true,         %
         breaklines=false,            % Zeilen werden Umgebrochen
         breakatwhitespace=true,     % Only allow breaks at white space
         keywordstyle=\textbf,
         frame=none,
 %        keywordstyle=[1]\textbf,    % Stil der Keywords
 %        keywordstyle=[2]\textbf,    %
 %        keywordstyle=[3]\textbf,    %
 %        keywordstyle=[4]\textbf,   \sqrt{\sqrt{}} %
 %        stringstyle=\color{white}\ttfamily, % Farbe der String
         showspaces=false,           % Leerzeichen anzeigen ?
         showtabs=false,             % Tabs anzeigen ?
         xleftmargin=34pt, %17pt,
         xrightmargin=34pt,
         framexleftmargin=17pt,
         framexrightmargin=5pt,
         framexbottommargin=4pt,
         captionpos=b,
         %backgroundcolor=\color{lightgray},
         showstringspaces=false,      % Leerzeichen in Strings anzeigeqn
                                % ?
         morekeywords={for,do,else,if,then,case,\?,ORIGIN,BODY,INCLUDE},%
         morecomment=[l]//,%
         string=[b]{'}
 }

\author{Thomas Bracht Laumann Jespersen\\ \texttt{laumann@cs.au.dk}}
\title{Learning \gbeta\\ A Tutorial\\\normalsize \gbeta version 1.9.8}

\begin{document}
\frontmatter

\maketitle

\begin{center}
  \includegraphics[scale=0.5]{pics/logo}
\end{center}

\begin{abstract}
  This is a book written to teach \gbeta . \gbeta is a generalization
  of \beta created by Erik Ernst at Aarhus University as the theme of
  his ph.d. thesis. \gbeta is implemented in \beta and as such, to
  understand \gbeta it is (somewhat) necessary to know \beta . This
  document aims to overcome this, and introduce the underlying
  concepts that \gbeta inherits directly from \beta , but without
  demanding any knowledge of \beta . In other words, the aim of this
  document is to teach \gbeta from the bottom. The only knowledge I
  will rely on is the understanding of object-oriented
  programming.

  To sum up, I seek to introduce the concepts of \gbeta along with
  actual examples to illustrate the (possible) applications of these
  concepts.

\end{abstract}

\newpage

\tableofcontents

\chapter{Motivation}

The motivation for writing this tutorial is mainly to combine the
learning material found at \url{http://www.daimi.au.dk/~eernst/gbeta},
Erik Ernst's own PhD thesis \cite{ernst99b} and the \beta book. In \gbeta version
1.9.8 the syntax has changed considerably, and \beta programs are
no longer valid \gbeta programs. The syntax change is partly to
conform to mainstream conventions (for instance using \verb|{ ... }| to denote
block structure) and to reduce the level of verbosity \beta (and
therefore alo \gbeta ) programs had a tendency to attain.

The outline of the book is as follows (subject to change):

\begin{description}\setlength{\itemsep}{-2pt}
\item[Assumptions] It's important to keep one's assumptions in mind. I
  assume the reader is familiar with object-oriented programming as a
  concept. Also I don't intend to explain how to install \gbeta . For
  information on how to obtain a working \gbeta compiler please refer
  to \url{http://www.daimi.au.dk/~eernst/gbeta}.
\item[Objects and Patterns] First we need to introduce the main concepts. The
  pattern will be introduced as the spear-heading abstraction
  introduced in \beta and preserved in \gbeta . Also the relation of
  patterns to object, how an object may be obtained from a pattern
  will be described
\item[Mixins] The concept of mixins should maybe be included. The
  question is: Is it relevant to the programmer? This determines
  whether or not it should be included.
\item[Virtual patterns] Introduce virtuals and their use.
\item[The fragment language] This can almost be included directly from
  Erik's website.
\end{description}

Note on the structure of the book: Consider dividing into two parts;
the first part should be the tutorial---all that is necessary for a
programmer to know to program in \gbeta . The second part should be
information aimed at the more language design curious reader. How are
things implemented and understood in \gbeta ?

\begin{flushright}
  Thomas Jespersen\\\today
\end{flushright}

\mainmatter

\chapter{\texttt{Hello, World!}}

\epigraph{I understand my tests are popular reading in the teachers' lounge}{Calvin}

No tutorial can rightfully call itself so if it lacks the
quintessential \texttt{Hello, World!} program. Accordingly, here is
the very first \gbeta program, \texttt{hello.gb}:

\lstinputlisting[caption=hello.gb]{gbsrc/hello.gb}

We can already note a few things from this example:

\begin{itemize}\setlength{\itemsep}{-2pt}
  \item Strings in \gbeta are a basic type. Literal strings are
    enclosed in single quotes (').
  \item The first line: \verb|-- universe:descriptor --|\ will be
    treated as `magic' for now, and explained in the section
    concerning the fragment system.
  \item The text \texttt{Hello, World!} is \emph{assigned} to
    \texttt{stdio}. In \gbeta\ assignments are on the form:
    \texttt{value|var}. Think of it as the value on the left side
    being `piped' into the right side.\footnote{bash programmers can
      think of the piping mechanism. For programmers used to the
      \texttt{var x = value} syntax this corresponds to \texttt{value|x}.}
\end{itemize}

\chapter{Objects and patterns}

This chapter introduces the reader to the \emph{pattern} concept, and
how objects work in \gbeta . Examples are given to illustrate the
usage of patterns and 

\begin{figure}[h!]
\caption{List of topics for this chapter}
\begin{itemize}\setlength{\itemsep}{-5pt}
\item the \textsf{MainPart} construct.
\item usage of the \textsf{MainPart} construct.
\item the \textsf{MainPart} is an example of a pattern.
\item static declarations, \verb|account1: @Account|.
\item dynamic declarations, \verb|account1: ^Account|.
\item declarations of the form \verb|account1: Account|.
\item the merge operator (\verb|&|).
\end{itemize}
\end{figure}

\section{The Pattern}

Assuming that the reader is familiar with the concept of an
object\footnote{and object-oriented programming} the first thing to
learn, understand and appreciate for any \gbeta\ programmer is the
concept of a \emph{pattern}. In \gbeta\ a pattern is declared like
this:
\begin{verbatim}
  Name : Merge
\end{verbatim}
Of course, this doesn't give us much information as we don't yet know
what \textsf{Merge} is. A pattern is the main abstraction used in
\gbeta , inherited from \beta . To understand what a pattern is and
how it works we must understand the concept of a
\textsf{MainPart}.\marginpar{The \textsf{MainPart} concept}

A \textsf{MainPart} is a piece of syntax that looks roughly like this:
\begin{verbatim}
  %( <Enter:Evaluation> | <Exit:Evaluation> ) {
    <Attributes>
  #
    <Statements>
  }
\end{verbatim}

\begin{itemize}\setlength{\itemsep}{-2pt}
\item The \texttt{Attributes} is a list of \emph{declarations}
\item The \texttt{Statements} section is also called the
  \emph{do-part}, and it's indicated by the do-part separator,
  \verb|#| .
\item The \texttt{Enter} evaluation specifies input properties,
  i.e. arguments accepted for assignment or procedure call usage. It
  is called the \emph{enter-part}.
\item The \texttt{Exit} evaluation (including the pipe) specifies the
  output properties, i.e. values delivered when evaluating the main
  part. This is (surprisingly) referred to as the \emph{exit-part}.
\item The \verb|%|, called \emph{public marker}, denotes whether or
  not this mainpart is public (as opposed to private).
\end{itemize}

\begin{quotation}\emph{
  All the elements are optional, so the list of declarations may be
  empty, and any selection of the enter-, do-, and exit-parts may be
  absent. This means that the main part is both a light-weight
  construct when only few of the basic elements are present, and it is
  a very rich construct when everything is there.\footnote{from
    \url{http://www.daimi.au.dk/~eernst/gbeta}}}
\end{quotation}

Say we wanted to model a bank account. Well, as a bank account isn't
an object that receives, does or outputs anything we don't need the
enter-, do- and exit-parts. It is simply an object that holds a
balance and on which we may perform withdrawals and deposits. To model
this in \gbeta\ we could write
\begin{verbatim}
  Account: %{
    balance: ... ;
    Withdraw: ... ;
    Deposit: ... ;
  }
\end{verbatim}
where \texttt{balance} represents the current balance of the account,
and \texttt{Withdraw} and \texttt{Deposit} represent the actions of
withdrawing and depositing money respectively. \texttt{balance},
\texttt{Withdraw} and \texttt{Deposit} are \emph{attributes} of the
\texttt{Account} pattern.

Using the \texttt{Account} pattern we may declare objects like this:
\begin{verbatim}
  myAccount: @Account;
  account1,account2: @Account;
\end{verbatim}
and these will be objects created from the \texttt{Account}
pattern. \texttt{myAccount}, \texttt{account1} and \texttt{account2}
can all be treated the same way, but they refer to different instances
of the \texttt{Account} pattern.

\subsection*{Basic patterns}
There are, as in any programming language, a set of pre-defined
entities that can be used. In \gbeta\ the basic patterns are
\begin{itemize}\setlength{\itemsep}{-2pt}
  \item \texttt{int}, pattern representing an integer
  \item \texttt{bool}, traditional boolean
  \item \texttt{string}, a basic type in \gbeta
  \item \texttt{float}, the only floating-point basic pattern
%  \item \texttt{char}
\end{itemize}

We could for instance use the \texttt{int} pattern when representing
the balance in our bank account:
\begin{verbatim}
  balance: @int;
\end{verbatim}

\subsection*{Declarations}
The syntactic element for describing attributes is a
\emph{declaration}. The following are all declarations:
\begin{verbatim}
  Account: %{ ... };
  myAccount: @Account;
  balance: @int;
  const: %(|'Hello');
  sqrt: %(i:int|i*i) { 'squaring a number'|stdio }
\end{verbatim}
In general, the \verb|:| syntax denotes a name being associated with
an entity.

\subsection*{Descriptions of actions}
The \texttt{Deposit} and \texttt{Withdraw} actions may now be
described as follows:
\begin{verbatim}
  Account: %{
    balance: @int;
    
    Deposit: %(amount:int|balance) {
      balance+amount|balance
    };
    
    Withdraw: %(amount:int|balance) {
      balance-amount|balance
    }
  }
\end{verbatim}
\texttt{Deposit} and \texttt{Withdraw} are patterns. The execution of
a deposit action requires the input \texttt{amount} (which is an
integer) and it exits with the value of balance. The usage of
\texttt{Deposit} may look like this:
\begin{verbatim}
  500|myAccount.Deposit|newBalance
\end{verbatim}
 where \texttt{myAccount} is an object instance of the
 \texttt{Account} pattern. The pipe (\verb$|$) is used in
 assignments. The execution of \texttt{Deposit} goes as follows: the
 value 500 is assigned to the \texttt{amount} attribute, then the
 imperative \texttt{balance+amount|balance} adds \texttt{amount} to
 \texttt{balance} and assigns the result to \texttt{balance}. Finally,
 the value of \texttt{balance} is exited.\footnote{Readers may think
   of the value exited as being \emph{returned}}. In this example the
 exited value is assigned to an object denoted by \texttt{newBalance}.

\subsection*{Imperatives}
The term \emph{imperatives} is used synonymously with
\emph{statements}. An imperative is a syntactic element for describing 
actions, of which the following are examples:
\begin{verbatim}
  balance+amount|balance;
  500|account1.Deposit|newBalance
\end{verbatim}

\subsection*{A complete program}
Until now, apart from the \texttt{Hello, World!} program, everything
presented has been parts of the \gbeta\ language. It is time to show
how these elements can be combined in our account modelling example:

\lstinputlisting[caption=account.gb]{gbsrc/account.gb}

The only comments in \gbeta\ are single line comments that begin with
\verb|//|. As the reader may notice the entire program follows the
form of a main part with an attributes section and a do-part
section---no enter- and exit-parts.

Execution then starts at the point marked \texttt{L1}, where all
accounts and integer variables have been declared. At \texttt{L2} the
variables \texttt{k1}, \texttt{k2} and \texttt{k3} hold values 100,
200 and 300 respectively. At \texttt{L3} program execution stops and
\texttt{k1} holds the value 170, \texttt{k2} holds the value 290 and
\texttt{k3} holds the value 210.

\subsection*{\texttt{Hello, World!} revisited}
We should now be at a stage where can rewrite our first program using
the pattern abstraction:

\lstinputlisting[caption=Hello\, World! with patterns]{gbsrc/hello_pattern.gb}

The important patterns to notice here are in the attributes section of
this main part. The \texttt{s: string} and \texttt{hello:
  \%(|'Hello')}. \texttt{s} is an object instance of the basic
\texttt{string} pattern and is used as such in the line
\verb$'World!\n'|s$.

\texttt{hello} is a little more interesting as it
is a pattern that consists solely of an exit-part! When evaluated
there is nothing to be entered, nothing to do (no do-part), only exit
(return) the value \verb|'Hello'|. A pattern that does nothing but
exiting a value, can be understood to model a constant.\footnote{Readers
  familiar with constants from other programming languages should
  appreciate the simplicity of this}

The last line
\begin{verbatim}
  hello+', '+s|stdio
\end{verbatim}
denotes the plus operator being applied to the evaluations of
\texttt{hello}, the string literal \verb|', '| and \texttt{s}. These
are all \texttt{string} types and the plus operator then has the
semantics of concatenating them, so the final result assigned
to \texttt{stdio} is (again, surprisingly) \verb|'Hello, World!\n'|.
\subsection*{A note on semicolons}
As the reader may have noticed, semicolons are used in an odd
manner. The reason is that sometimes they are optional, sometimes they
are not. The general rule is the that \emph{the last statement in a
  list of statements need not be followed by a semicolon}. In other words
semicolons are mandatory to separate statemements (imperatives and
declarations), but at the end of a construct (pattern, control
structure, etc.) the semicolon is optional. Why this is so, is a
discussion I will omit at this stage. If one wishes to avoid compiler
complaints due to lacking semicolons, just always insert it.

\section{On Declarations}

This section discusses different kinds of declarations --- especially
the difference static and dynamic declarations.

\subsection*{Static and dynamic declarations}

We have already seen static declarations, they are on the form
\verb|account1: @Account|. A static declaration like this implies
creation of an object from the given pattern (in this case
\texttt{Account}), and execution of its do-part.

The same can be said for dynamic declarations, which look like this:
\verb|account1: ^Account|, the difference being that \texttt{account1}
now refers to the \emph{address} of the \texttt{Account} object, where
the name of the statically declared variable refers to the object
itself (no addresses involved).

\chapter{Control Structures}

% Chapter about the various control structures of gbeta

\begin{itemize}\setlength{\itemsep}{-2pt}
\item \texttt{if} statements
\item \texttt{for} statements
\item \texttt{case} statements
\item \texttt{label},\texttt{leave} and \texttt{restart} statements
\end{itemize}

\section{The \texttt{if} statement}

An \texttt{if} statement is a conditional statement, where a program's
behaviour can be switched based on a boolean expression. An \texttt{if}
statement can look in one of the two following ways:

\begin{verbatim}
  if <Evaluation> do { <Statements> }
\end{verbatim}
or:
\begin{verbatim}
  if <Evaluation> do { <Statements> }
  else { <Statements> }
\end{verbatim}


%%% Local Variables: 
%%% mode: latex
%%% TeX-master: "gbeta-tutorial"
%%% End: 


\chapter{Virtual Patterns}

In this chapter the reader is introduced to the \emph{virtual
  pattern}, and \emph{virtual declarations}. More about this later...

\subsection*{Virtual pattern declaration}

Maybe we need some text here\dots

\newpage

The program \texttt{99 bottles of beer} in \gbeta :

\lstinputlisting[caption=99 bottles of beer]{gbsrc/beer.gb}

\bibliographystyle{plain}
\bibliography{gbeta-tutorial}

\end{document}
